\chapter[The symmetric group and Specht modules]{The Symmetric Group and Specht Modules}

Write \(S_n\) for the permutations of \(\{1,\ldots,n\}\), composed from
right to left. All modules in the main argument are complex and
finite-dimensional. We use the group algebra \(A_n=\mathbb C[S_n]\) and
regard its representations as left modules. The sign of \(g\in S_n\)
is written \(\operatorname{sgn}(g)\). The group \(S_0\) is trivial, and
every construction on zero labels uses the empty diagram.

The finite-group results used below are Schur's lemma, Maschke's
theorem and the count of irreducible complex representations by
conjugacy classes. They are established in \cite[Chapter 8]{algebra-core}.
For the count, the finite-group theorem starts from a complete family
of irreducibles. Its characters are an orthonormal basis of the space
of class functions, whose dimension is the number of conjugacy classes.
We apply this count to determine when a constructed family is complete.

\section{Cycle types and permutation modules}

\begin{theorem}[Conjugacy classes and their sizes]
\label{res:A01}
The cycle type of a permutation records all its cycle lengths, including
one for each fixed point. It gives a bijection from conjugacy classes of
\(S_n\) to partitions of \(n\). If \(\nu\) has \(m_j\) parts of size \(j\),
the centralizer of a permutation of type \(\nu\) has order
\[
  z_\nu=\prod_{j\geq 1}j^{m_j}m_j!,
\]
and the corresponding class has \(n!/z_\nu\) elements. At \(n=0\) these
counts are both \(1\).
\end{theorem}

\begin{proof}
Conjugation relabels each orbit of a permutation and preserves its
cyclic order. If two permutations have the same type, pair their cycles
of equal length and choose a starting label in each paired cycle.
Mapping successive labels in the source cycle to successive labels in
its target defines a bijection that conjugates the permutations.
Conversely, for any partition divide the labels into successive blocks of
its part sizes and cycle the labels of each block. This realizes every
partition, with the sole possibility at \(n=0\) being the identity.

To commute with a given permutation \(g\), a permutation must send each
cycle of \(g\) to a cycle of the same length. For cycles of length \(j\)
there are \(m_j!\) choices of their target order and \(j\) choices of
the image of one chosen starting label on each cycle. Commuting with
\(g\) then determines the images of all successive labels. Each such
choice gives a bijection commuting with \(g\), independently for each
length. The class-size formula now follows from orbit--stabilizer for
conjugation. See \cite[\S2.1]{sagan} and \cite[\S4.14]{etingof}.
\end{proof}

A partition \(\lambda\vdash n\) is a weakly decreasing list of positive
integers summing to \(n\). We may append zeros when comparing lists.
Its Young diagram consists of cells \((i,j)\) with \(1\leq j\leq\lambda_i\).
Rows are numbered from top to bottom and columns from left to right.
The transpose partition \(\lambda^\top\) has \(j\)-th part equal to the
height of column \(j\), and \(\ell(\lambda)\) is the number of nonempty rows.

\begin{proposition}[Tableaux, tabloids, and induction]
\label{res:A02}
A \(\lambda\)-tableau \(t\) is a bijective labeling of the diagram of
\(\lambda\) by \(\{1,\ldots,n\}\). Two tableaux give the same tabloid if
the sets of labels in the corresponding rows agree. Write \(\{t\}\)
for the tabloid of \(t\). The row subgroup \(R_t\) and the column
subgroup \(C_t\) permute labels within the rows and columns of \(t\),
respectively, so \(R_t\cap C_t=\{1\}\). The left permutation module
\(M^\lambda\) on the basis of tabloids is isomorphic to
\(\operatorname{Ind}_{R_t}^{S_n}\mathbf 1\), where \(\mathbf 1\) is the
trivial \(R_t\)-module.
\end{proposition}

\begin{proof}
For a fixed tableau \(t\), any other tableau is \(gt\) for a unique
permutation \(g\), where \(g\) replaces each entry \(a\) by \(g(a)\).
The stabilizer of \(\{t\}\) is exactly \(R_t\) because the tabloid
remembers the set of labels in each row. A permutation preserving both
row sets and column sets fixes the label in their intersection at each
cell. Thus \(R_t\cap C_t\) is trivial. Also \(R_t\) is a product of
symmetric groups on the row sets, the Young subgroup for this choice
of labeling.

With the induction model from \cite[Chapter 9]{algebra-core}, send
\(g\otimes 1\) in \(A_n\otimes_{\mathbb C[R_t]}\mathbb C\) to \(\{gt\}\).
The relation \(gr\otimes 1=g\otimes 1\) holds for every \(r\in R_t\),
and the proposed map respects it. As a right \(\mathbb C[R_t]\)-module,
\(A_n\) is a direct sum of the spans \(g\mathbb C[R_t]\), one for each
left coset \(gR_t\). Tensoring makes each summand a line and sends its
generator to the corresponding tabloid. Hence the map is bijective
and intertwines left multiplication by \(S_n\). See
\cite[\S2.1, Theorem 2.1.12]{sagan}.
\end{proof}

For partitions \(\lambda,\mu\vdash n\), write \(\lambda\unrhd\mu\) if
\(\sum_{i=1}^k\lambda_i\geq\sum_{i=1}^k\mu_i\) for every \(k\).
These sums include appended zeros. In particular, dominance in both
directions forces equal parts.

\begin{lemma}[The row--column test]
\label{res:A03}
If \(\lambda\not\unrhd\mu\), then for any \(\lambda\)-tableau \(t\) and
\(\mu\)-tableau \(u\), some column of \(t\) and some row of \(u\) have at
least two labels in common. If \(\lambda=\mu\) and there is no such
intersection, then \(\{u\}=\{ct\}\) for a unique \(c\in C_t\).
\end{lemma}

\begin{proof}
In the absence of an intersection of size at least two, column \(j\)
of \(t\) has at most \(\min(k,\lambda^\top_j)\) of its labels in the
first \(k\) rows of \(u\). Counting all the labels in those rows gives
\[
  \sum_{i=1}^k\mu_i
  \ \leq\ \sum_{j\geq 1}\min(k,\lambda^\top_j)
  \ =\ \sum_{i=1}^k\lambda_i .
\]
The last equality counts the cells in the first \(k\) rows of the
\(\lambda\)-diagram by columns. This proves the first claim. When the
shapes are equal, the sum is an equality for each \(k\), so every
column attains its individual upper bound. Taking the differences of
these bounds at consecutive \(k\) shows that each column \(j\) supplies
exactly one label to row \(i\) of \(u\) if \(i\leq\lambda^\top_j\),
and none otherwise. Within each column, send the label of \(t\) in
cell \((i,j)\) to that column's unique label in row \(i\) of \(u\).
The resulting \(c\in C_t\) has \(\{ct\}=\{u\}\). Two such choices
could differ only by an element of \(C_t\cap R_t\), which is trivial.
Compare \cite[\S\S2.2--2.3]{sagan}.
\end{proof}

\section{Polytabloids and irreducibility}

\begin{proposition}[Column alternation and the Specht module]
\label{res:A04}
For a tableau \(t\), put
\[
  b_t=\sum_{c\in C_t}\operatorname{sgn}(c)c\in A_n,
  \qquad e_t=b_t\{t\}\in M^\lambda .
\]
The vector \(e_t\) is called a polytabloid. Its coefficient at \(\{t\}\)
is \(1\), and \(g e_t=e_{gt}\) for all \(g\in S_n\).
The Specht module \(S^\lambda\) is the span of the polytabloids of shape
\(\lambda\). It is a nonzero submodule generated by any one \(e_t\).
For \(c\in C_t\), \(e_{ct}=\operatorname{sgn}(c)e_t\), and
\(b_t e_t=|C_t|e_t\).
\end{proposition}

\begin{proof}
By Proposition~\ref{res:A02}, only the identity in \(C_t\) fixes
\(\{t\}\); in fact the \(\{ct\}\) for \(c\in C_t\) are distinct.
Since \(C_{gt}=gC_tg^{-1}\) and sign is invariant under conjugation,
\(b_{gt}=gb_tg^{-1}\), which gives \(ge_t=e_{gt}\). The action on
tableaux is transitive, so the orbit of any \(e_t\) generates the
indicated nonzero module. Multiplying \(b_t\) on either side by
\(c\in C_t\) gives \(\operatorname{sgn}(c)b_t\) by reindexing the sum.
Consequently \(c e_t=\operatorname{sgn}(c)e_t\), and each signed
summand of \(b_t\) acts on \(e_t\) as the identity. This yields
\(b_t e_t=|C_t|e_t\). See \cite[\S2.3]{sagan}.
\end{proof}

Give \(M^\lambda\) the positive definite Hermitian inner product
with orthonormal tabloid basis, linear in the first variable. The
permutation action preserves this inner product.

\begin{theorem}[Submodule theorem]
\label{res:A05}
For a tableau \(t\) of shape \(\lambda\), the action of \(b_t\) on
\(M^\lambda\) is \(b_t x=\langle x,e_t\rangle e_t\).
Every submodule \(U\subseteq M^\lambda\) contains \(S^\lambda\) or
is contained in \((S^\lambda)^\perp\). In particular \(S^\lambda\)
is irreducible.
\end{theorem}

\begin{proof}
First consider the action on an arbitrary tabloid \(\{u\}\), possibly
of a different shape. If some row of \(u\) contains two labels in a
column of \(t\), the transposition of these labels belongs to \(C_t\)
and fixes \(\{u\}\). Pairing terms in \(b_t\) by right multiplication
by this transposition gives \(b_t\{u\}=0\). In the equal-shape case
with no such collision, Lemma~\ref{res:A03} gives
\(\{u\}=c\{t\}\) for \(c\in C_t\). The product rule for \(b_t c\)
then gives \(b_t\{u\}=\operatorname{sgn}(c)e_t\).

Thus the image on \(M^\lambda\) is the line \(\mathbb C e_t\). The
adjoint of each basis permutation is its inverse, and
\(\operatorname{sgn}(c^{-1})=\operatorname{sgn}(c)\), so \(b_t\)
is self-adjoint. If \(b_t x=\alpha e_t\), taking the inner product
against \(\{t\}\) gives
\[
  \alpha=\langle b_t x,\{t\}\rangle
        =\langle x,b_t\{t\}\rangle=\langle x,e_t\rangle .
\]
Here the first equality uses the coefficient \(1\) from
Proposition~\ref{res:A04}. If some vector \(x\in U\) pairs nontrivially
with any \(e_t\), applying \(b_t\) puts that \(e_t\) in \(U\), and the
orbit of \(e_t\) gives \(S^\lambda\subseteq U\). Otherwise \(U\)
is perpendicular to every \(e_t\) and hence to \(S^\lambda\).
Apply this dichotomy to a submodule of \(S^\lambda\). In the second
case it is zero by positive definiteness, while in the first it is all
of \(S^\lambda\). This is the argument of \cite[\S2.4]{sagan}.
\end{proof}

\begin{proposition}[Dominance and diagonal multiplicity]
\label{res:A06}
If \(\operatorname{Hom}_{S_n}(S^\lambda,M^\mu)\neq 0\), then
\(\lambda\unrhd\mu\). The space
\(\operatorname{Hom}_{S_n}(S^\lambda,M^\lambda)\) has dimension \(1\).
\end{proposition}

\begin{proof}
An intertwiner \(f:S^\lambda\to M^\mu\) is determined by \(f(e_t)\)
for a fixed \(\lambda\)-tableau \(t\), since \(e_t\) generates.
Moreover
\[
  b_t f(e_t)=f(b_t e_t)=|C_t| f(e_t),
\]
so \(f(e_t)\) belongs to the image of \(b_t\). If
\(\lambda\not\unrhd\mu\), the first part of Lemma~\ref{res:A03}
and the transposition pairing in Theorem~\ref{res:A05} show
that \(b_t\) kills every basis tabloid of \(M^\mu\).
If \(\lambda=\mu\), Theorem~\ref{res:A05} gives the one-dimensional
image \(\mathbb C e_t\), and the inclusion \(S^\lambda\hookrightarrow
M^\lambda\) gives a nonzero intertwiner. See
\cite[\S2.4, Proposition 2.4.5]{sagan}.
\end{proof}

\begin{theorem}[Classification over \(\mathbb C\)]
\label{res:A07}
The modules \(S^\lambda\), for \(\lambda\vdash n\), form a complete
family of pairwise nonisomorphic irreducible \(S_n\)-modules.
For \(n\geq 1\), \(S^{(n)}\) is trivial and \(S^{(1^n)}\) is the
sign module. At \(n=0\), \(S^\varnothing\) is the trivial line.
\end{theorem}

\begin{proof}
Theorem~\ref{res:A05} proves irreducibility. If
\(S^\lambda\simeq S^\mu\), then composing an isomorphism and its
inverse with the respective inclusions in \(M^\mu\) and \(M^\lambda\)
gives nonzero intertwiners. Proposition~\ref{res:A06} yields both
dominance inequalities, so \(\lambda=\mu\).
Theorem~\ref{res:A01} gives a conjugacy class for each partition,
so the number constructed is already the number of isomorphism
classes given by the finite-group count.

For one row there is one tabloid and every permutation fixes it.
For one column, \(C_t=S_n\), and Proposition~\ref{res:A04} makes
the orbit of \(e_t\) the sign line. The empty diagram has one tabloid
and the only permutation fixes it. The classification is
\cite[\S2.4, Theorem 2.4.6]{sagan}; compare
\cite[\S4.12, Theorem 4.36]{etingof}.
\end{proof}

\section{Young operators in the group algebra}

\begin{proposition}[Left ideals and normalized Young idempotents]
\label{res:A08}
Fix a tableau \(t\) of shape \(\lambda\) and put
\[
  a_t=\sum_{r\in R_t}r,\qquad c_t=b_t a_t,\qquad c'_t=a_t b_t .
\]
The map \(\{gt\}\mapsto ga_t\) identifies \(M^\lambda\) with
\(A_n a_t\), and its restriction identifies \(S^\lambda\) with
\(A_n c_t\). With \(q_\lambda=n!/\dim S^\lambda\),
\[
  c_t^2=q_\lambda c_t,\qquad (c'_t)^2=q_\lambda c'_t .
\]
Thus \(c_t/q_\lambda\) and \(c'_t/q_\lambda\) are idempotents, and
the left ideals \(A_n c_t\) and \(A_n c'_t\) are isomorphic.
\end{proposition}

\begin{proof}
The map on tabloids is well-defined since \(ra_t=a_t\) for \(r\in R_t\).
The sum \(ga_t\) has support on \(gR_t\), and sums on distinct cosets
have disjoint supports. Hence the map identifies the tabloid basis
with a basis of \(A_n a_t\) and is equivariant. It sends \(e_t\)
to \(b_t a_t=c_t\), which generates the image \(A_n c_t\).

Right multiplication by \(c_t\) on \(A_n\) is a left-module map with
image \(A_n c_t\). Its restriction to that image is multiplication
by some scalar \(q\) by Schur's lemma, since \(A_n c_t\simeq S^\lambda\)
is irreducible. In the group basis its trace is \(n!\) times the
coefficient of \(1\) in \(c_t\). That coefficient is \(1\): an equality
\(cr=1\) for \(c\in C_t\), \(r\in R_t\) implies \(c=r=1\).
The map induces zero on the quotient \(A_n/(A_n c_t)\), so its trace
also equals \(q\dim S^\lambda\). Thus \(q=q_\lambda\), and applying
the restriction to \(c_t\) gives \(c_t^2=q_\lambda c_t\).
The linear map \(g\mapsto g^{-1}\) reverses products and fixes \(a_t\)
and \(b_t\), so it takes the square relation for \(c_t\) to the
one for \(c'_t\).

Right multiplication by \(b_t\) maps \(A_n c_t\) into \(A_n c'_t\)
because \(c_t b_t=b_t c'_t\). Right multiplication by \(a_t\) maps
back because \(c'_t a_t=a_t c_t\). The compositions are right
multiplication by \(c_t\) and \(c'_t\) on the respective ideals,
each equal to scalar \(q_\lambda\) by the square relations. These
maps are therefore inverse up to a nonzero scalar. Our column-sum
then row-sum order is the \(b_t a_t\) model compared with \(a_t b_t\)
in \cite[Lemma 4.42 and Problem 4.68(a)]{etingof}.
\end{proof}

\begin{corollary}[Transpose and sign twist]
\label{res:A09}
For every \(\lambda\vdash n\),
\(S^{\lambda^\top}\simeq S^\lambda\otimes\operatorname{sgn}\).
The regular module of \(S_2\) over \(\mathbb F_2\) has a fixed line
with no invariant complement. It illustrates the role of invertibility
of the group order in Maschke's theorem.
\end{corollary}

\begin{proof}
The algebra automorphism \(\theta\) defined by
\(\theta(g)=\operatorname{sgn}(g)g\) is an involution. For a left
ideal \(I\), the linear map \(x\mapsto\theta(x)\) intertwines the
sign-twisted action on \(I\) with the ordinary action on \(\theta(I)\),
since \(\theta(\operatorname{sgn}(g)gx)=g\theta(x)\).
The transposed tableau \(t^\top\) interchanges the row and column
subgroups. Hence \(\theta(b_t)=a_{t^\top}\),
\(\theta(a_t)=b_{t^\top}\), and \(\theta(c_t)=c'_{t^\top}\).
Proposition~\ref{res:A08} identifies \(A_n c'_{t^\top}\) with
\(S^{\lambda^\top}\), proving the sign twist formula.

For \(S_2=\{1,s\}\) over \(\mathbb F_2\), the regular module has
basis \(1,s\), and \(s\) exchanges the two basis vectors. The fixed
space is the line spanned by \(1+s\). An invariant complement would
have dimension \(1\), and on it \(s\) would act by a nonzero scalar
in \(\mathbb F_2\), necessarily \(1\). Such a complement would
therefore lie in the fixed line, which is impossible.
Compare \cite[\S3.1 and Problem 4.68(b)]{etingof}.
\end{proof}
